\documentclass[answers,12pt,addpoints]{exam} 
\usepackage{import}

\import{C:/Users/prana/OneDrive/Desktop/MathNotes}{style.tex}

% Header
\newcommand{\name}{Pranav Tikkawar}
\newcommand{\course}{01:XXX:XXX}
\newcommand{\assignment}{Homework n}
\author{\name}
\title{\course \ - \assignment}

\begin{document}
\maketitle
\tableofcontents
\newpage
\section{L1}

\section{L2}

\section{L3}

\section{L4}

\begin{definition}[Conditional Expectation]
Let $(\Omega,\mathcal{F},\mathbb{P})$ be a probability space, $X$ an integrable random variable with $\mathbb{E}|X|<\infty$, and $\mathcal{G}\subseteq\mathcal{F}$ a sub-$\sigma$-algebra. A random variable $Y$ is called the \emph{conditional expectation} of $X$ given $\mathcal{G}$, written
\[
Y = \mathbb{E}[X \mid \mathcal{G}],
\]
if
\begin{enumerate}
  \item $Y$ is $\mathcal{G}$-measurable, and
  \item for every $A\in\mathcal{G}$, $\displaystyle \int_A Y\,d\mathbb{P} = \int_A X\,d\mathbb{P}$.
\end{enumerate}
When $Z$ is a random variable, we write $\mathbb{E}[X\mid Z] := \mathbb{E}[X\mid \sigma(Z)]$.
\end{definition}

\section{L5}

\begin{definition}[Filtration]
A \emph{filtration} on a probability space $(\Omega,\mathcal{F},\mathbb{P})$ is an increasing sequence of sub-$\sigma$-algebras
\[
\mathcal{F}_0 \subseteq \mathcal{F}_1 \subseteq \mathcal{F}_2 \subseteq \cdots \subseteq \mathcal{F}.
\]
We interpret $\mathcal{F}_n$ as the information available up to time $n$.
\end{definition}

\begin{definition}[Martingale]
Let $(\mathcal{F}_n)_{n\ge 0}$ be a filtration. A stochastic process $(Z_n)_{n\ge 0}$ is a \emph{martingale} with respect to $(\mathcal{F}_n)$ if:
\begin{enumerate}
  \item $Z_n$ is $\mathcal{F}_n$-measurable for each $n$ (adapted),
  \item $\mathbb{E}|Z_n|<\infty$ for each $n$, and
  \item $\mathbb{E}[Z_{n+1}\mid \mathcal{F}_n] = Z_n$ almost surely for all $n\ge 0$.
\end{enumerate}
\end{definition}

\begin{remark}[Martingale, Submartingale, Supermartingale]
With the same setup:
\begin{itemize}
  \item $(Z_n)$ is a \emph{submartingale} if $\mathbb{E}[Z_{n+1}\mid \mathcal{F}_n] \ge Z_n$ a.s.\ for all $n$.
  \item $(Z_n)$ is a \emph{supermartingale} if $\mathbb{E}[Z_{n+1}\mid \mathcal{F}_n] \le Z_n$ a.s.\ for all $n$.
\end{itemize}
If $(Z_n)$ is a martingale, then $\mathbb{E}Z_n = \mathbb{E}Z_0$ for all $n$.
\end{remark}

\begin{definition}[Stopping Time]
Given a filtration $(\mathcal{F}_n)_{n\ge 0}$, a random variable $T:\Omega\to \{0,1,2,\dots\}\cup\{\infty\}$ is a \emph{stopping time} if
\[
\{T = n\} \in \mathcal{F}_n \quad \text{for every } n.
\]
Equivalently, $\{T\le n\}\in\mathcal{F}_n$ for each $n$.
\end{definition}

\begin{theorem}[Wald's Lemma I / Optional Stopping, Bounded Case]
Let $(\mathcal{F}_n)_{n\ge 0}$ be a filtration and $(Z_n)_{n\ge 0}$ a martingale with respect to $(\mathcal{F}_n)$. Let $T$ be a stopping time such that $T\le N$ almost surely for some deterministic $N$. Then
\[
\mathbb{E}[Z_T] = \mathbb{E}[Z_0].
\]
\end{theorem}

\section{L6}

\begin{definition}[Almost Sure Convergence]
We say $X_n \to X$ \emph{almost surely} (a.s.) if
\[
\mathbb{P}\big(\{\omega\in\Omega : X_n(\omega)\to X(\omega)\}\big) = 1.
\]
\end{definition}

\begin{definition}[Convergence in Expectation]
We say $X_n \to X$ \emph{in expectation} if $\mathbb{E}|X_n - X|\to 0$. (Often called convergence in $L^1$.)
\end{definition}

\begin{definition}[Monotone Convergence Theorem (MCT)]
Let $(X_n)_{n\ge 1}$ be nonnegative random variables with $0\le X_1 \le X_2 \le \cdots$ and $X_n \to X$ a.s. Then
\[
\lim_{n\to\infty}\mathbb{E}[X_n] = \mathbb{E}[X].
\]
\end{definition}

\begin{definition}[Dominated Convergence Theorem (DCT)]
Suppose $X_n \to X$ almost surely, and there exists an integrable random variable $Y$ such that $|X_n|\le Y$ a.s.\ for all $n$ and $\mathbb{E}|Y|<\infty$. Then
\[
\lim_{n\to\infty}\mathbb{E}[X_n] = \mathbb{E}[X].
\]
\end{definition}

\section{L7}

\begin{lemma}[Convex Transform of a Submartingale]
Let $(Z_n)$ be a submartingale with respect to $(\mathcal{F}_n)$ and let $\varphi:\mathbb{R}\to\mathbb{R}$ be convex and nondecreasing. Assume $\mathbb{E}|\varphi(Z_n)|<\infty$ for every $n$. Then $(\varphi(Z_n))$ is a submartingale.
\end{lemma}

\begin{theorem}[Not-To-Gamble Theorem]
Let $(X_n)_{n\ge 0}$ be a supermartingale with respect to $(\mathcal{F}_n)$, and let $(H_n)_{n\ge 1}$ be a predictable sequence (i.e.\ $H_n$ is $\mathcal{F}_{n-1}$-measurable) with $H_n\ge 0$ a.s. Define
\[
(H\cdot X)_n := \sum_{i=1}^n H_i (X_i - X_{i-1}), \quad (H\cdot X)_0 := 0.
\]
Then $((H\cdot X)_n)$ is a supermartingale. In particular, no nonnegative predictable betting strategy can turn a losing game (supermartingale) into a winning game in expectation.
\end{theorem}

\begin{theorem}[Optional Stopping, Bounded Case for Sub/Supermartingales]
Let $(X_n)_{n\ge 0}$ be adapted to $(\mathcal{F}_n)$ and let $T$ be a stopping time with $T\le N$ a.s. If $(X_n)$ is a submartingale, then
\[
\mathbb{E}[X_N] \ge \mathbb{E}[X_T] \ge \mathbb{E}[X_0].
\]
If $(X_n)$ is a supermartingale, then
\[
\mathbb{E}[X_N] \le \mathbb{E}[X_T] \le \mathbb{E}[X_0].
\]
\end{theorem}

\section{L8}

\begin{theorem}[Submartingale and Supermartingale Convergence]
Let $(X_n)_{n\ge 0}$ be a submartingale. If
\[
\sup_{n\ge 0} \mathbb{E}[X_n^{+}] < \infty,
\]
then $X_n$ converges almost surely to some integrable random variable $X_\infty$ and $\mathbb{E}|X_\infty|<\infty$.

Let $(Z_n)_{n\ge 0}$ be a supermartingale. If there exists $c\in\mathbb{R}$ such that $Z_n \ge c$ a.s.\ for all $n$, then $Z_n$ converges almost surely to some integrable random variable $Z$ and $\mathbb{E}[Z]\le \mathbb{E}[Z_0]$.
\end{theorem}


\section{L9}

\section{L10}

\begin{theorem}[Markov Inequality]
Let $X$ be a nonnegative random variable and $c>0$. Then
\[
\mathbb{P}(X \ge c) \le \frac{\mathbb{E}[X]}{c}.
\]
\end{theorem}

\begin{theorem}[Chebyshev Inequality]
Let $X$ be a random variable with $\mathbb{E}X<\infty$ and $\mathrm{Var}(X)<\infty$. Then for any $c>0$,
\[
\mathbb{P}\big(|X-\mathbb{E}X|\ge c\big) \le \frac{\mathrm{Var}(X)}{c^2}.
\]
\end{theorem}

\begin{theorem}[Kolmogorov Maximal Inequality]
Let $X_1,\dots,X_n$ be independent random variables with $\mathbb{E}X_i=0$ and $\mathrm{Var}(X_i)<\infty$. Define partial sums $S_k:=\sum_{i=1}^k X_i$ for $k=1,\dots,n$, and $S_0=0$. Then for any $\lambda>0$,
\[
\mathbb{P}\Big(\max_{1\le k\le n} |S_k|\ge \lambda\Big) \le \frac{\mathrm{Var}(S_n)}{\lambda^2}
= \frac{\sum_{i=1}^n \mathrm{Var}(X_i)}{\lambda^2}.
\]
\end{theorem}

\begin{theorem}[Doob's Maximal Inequality]
Let $(M_k)_{0\le k\le n}$ be a submartingale. Then for any $\lambda>0$,
\[
\mathbb{P}\Big(\max_{0\le k\le n} M_k \ge \lambda\Big) \le \frac{\mathbb{E}[M_n^{+}]}{\lambda},
\]
where $M_n^{+} = \max\{M_n,0\}$.
\end{theorem}

\begin{theorem}[Ville's Inequality]
Let $(L_k)_{k\ge 0}$ be a nonnegative supermartingale with respect to $(\mathcal{F}_k)$. Then for any $x>0$,
\[
\mathbb{P}\Big(\sup_{k\ge 0} L_k \ge x\Big) \le \frac{\mathbb{E}[L_0]}{x}.
\]
\end{theorem}

\begin{theorem}[Fatou's Lemma]
Let $(X_n)_{n\ge 1}$ be nonnegative random variables. Then
\[
\mathbb{E}\Big[\liminf_{n\to\infty} X_n\Big] \le \liminf_{n\to\infty} \mathbb{E}[X_n].
\]
\end{theorem}

\section{L11}

\begin{definition}[$\sigma^2$-Subgaussian Random Variable]
A random variable $X$ is called \emph{$\sigma^2$-subgaussian} if $\mathbb{E}X=0$ and
\[
\mathbb{E}[e^{sX}] \le \exp\left(\frac{\sigma^2 s^2}{2}\right)\quad \text{for all } s\in\mathbb{R}.
\]
\end{definition}

\begin{lemma}[Hoeffding's Lemma and Hoeffding's Inequality]
If $X$ satisfies $a\le X\le b$ a.s.\ and $\mathbb{E}X=0$, then for all $s\in\mathbb{R}$,
\[
\mathbb{E}[e^{sX}] \le \exp\left(\frac{s^2(b-a)^2}{8}\right),
\]
so $X$ is $(b-a)^2/4$-subgaussian.

If $X_1,\dots,X_n$ are independent random variables with $a_i \le X_i \le b_i$ a.s., and $\bar X_n := \frac{1}{n}\sum_{i=1}^n X_i$, then for any $t>0$,
\[
\mathbb{P}\big(\bar X_n - \mathbb{E}\bar X_n \ge t\big) \le \exp\left(-\frac{2n^2 t^2}{\sum_{i=1}^n (b_i - a_i)^2}\right),
\]
and similarly for the two-sided deviation.
\end{lemma}

\begin{theorem}[Azuma--Hoeffding Inequality for Martingales]
Let $(Z_n)_{n\ge 0}$ be an integrable martingale with respect to $(\mathcal{F}_n)$, and assume there exist constants $\alpha_i,\beta_i\ge 0$ such that
\[
-\alpha_i \le Z_i - Z_{i-1} \le \beta_i \quad \text{a.s. for all } i\ge 1.
\]
Then for any $n\ge 1$ and $t>0$,
\[
\mathbb{P}(Z_n - \mathbb{E}Z_n \ge t) \le \exp\left(-\frac{2t^2}{\sum_{i=1}^n (\alpha_i + \beta_i)^2}\right),
\]
and the same bound holds for $\mathbb{P}(|Z_n - \mathbb{E}Z_n|\ge t)$ up to a factor 2.
\end{theorem}

\begin{theorem}[McDiarmid's Inequality]
Let $Z = h(X_1,\dots,X_n)$ where $X_1,\dots,X_n$ are independent, and suppose $h$ satisfies the bounded difference condition: whenever $x$ and $\tilde x$ differ only in coordinate $i$, we have
\[
|h(x) - h(\tilde x)| \le s.
\]
Then for any $t>0$,
\[
\mathbb{P}(Z - \mathbb{E}Z \ge t) \le \exp\left(-\frac{2t^2}{s^2 n}\right),
\]
and similarly for the two-sided deviation with an extra factor 2.
\end{theorem}

\section{L12}

\begin{definition}[Transition Probability Matrix]
For a time-homogeneous Markov chain $(X_n)$ on a finite state space $S=\{x_1,\dots,x_m\}$, the \emph{transition matrix} $P\in\mathbb{R}^{m\times m}$ is defined by
\[
P(i,j) := \mathbb{P}(X_{n+1} = x_j \mid X_n = x_i),
\]
which is independent of $n$. Each row of $P$ is a probability vector.
\end{definition}

\begin{example}
  [Biased Random Walk with absorbing Boundary]
  $$ h(i) =\begin{cases}
    \frac{1 - (q/p)^i}{1 - (q/p)^N} & \text{if } p\neq q,\\
    \frac{i}{N} & \text{if } p=q,
  \end{cases} $$
  where $h(i) = \mathbb{P}(\text{hit } N \text{ before } 0 \mid X_0 = i)$.
\end{example}

\begin{theorem}[Chapman--Kolmogorov Equation (Finite Case)]
For a time-homogeneous Markov chain with transition matrix $P$, we have
\[
\mathbb{P}(X_{n+k} = x_j \mid X_n = x_i) = (P^k)(i,j)
\]
for all $n,k\ge 0$ and all states $x_i,x_j\in S$.
\end{theorem}

\begin{definition}[Stationary Distribution]
A probability vector $\pi$ on $S$ is called a \emph{stationary distribution} for a Markov chain with transition matrix $P$ if
\[
\pi = \pi P.
\]
Equivalently, if $X_0\sim\pi$, then $X_n\sim\pi$ for all $n\ge 0$.
\end{definition}

\begin{definition}[Detailed Balance and Reversible Chains]
A distribution $\pi$ on $S$ satisfies the \emph{detailed balance} equations with respect to $P$ if
\[
\pi(i)P(i,j) = \pi(j)P(j,i)\quad \text{for all } i,j.
\]
In this case the chain is called \emph{reversible} with respect to $\pi$, and $\pi$ is automatically stationary.
\end{definition}

\begin{definition}
  [Metropolis--Hastings Algorithm]
  fix a distribution $q(x,\cdot)$ on $S$ for each $x\in S$. Given a target distribution $\pi$ on $S$, we define a Markov chain with transition matrix.
\[
P(x,y) = q(x,y) \alpha(x,y) \quad \text{for } y\neq x, \quad P(x,x) = 1 - \sum_{y\neq x} P(x,y),
\]
where
\[
\alpha(x,y) = \min\left(1, \frac{\pi(y)q(y,x)}{\pi(x)q(x,y)}\right).
\]
Set $X_{n+1} = \begin{cases}
y & \text{with probability } P(X_n,y),\\
X_n & \text{with probability } P(X_n,X_n).
\end{cases}$
Then $\pi$ is a stationary distribution for this chain.
\end{definition}


\begin{definition}
  [Gibbs Sampler]
  Let $\pi$ be a distribution on $S = S_1 \times \cdots \times S_d$. For each $i$, define the conditional distribution
\[
\pi_i(x_i \mid x_{-i}) := \mathbb{P}_{X\sim\pi}(X_i = x_i \mid X_j = x_j \text{ for } j\neq i).
\]
The Gibbs sampler is a Markov chain on $S$ that updates one coordinate at a time
in a random ordeer. We choose $i$ uniformly at random from $\{1,\dots,d\}$, and then update $X_i$ according to $\pi_i(\cdot \mid X_{-i})$. This chain has $\pi$ as a stationary distribution.   
\end{definition}



\section{L13}

\section{L14}

\begin{remark}[Extended Stationary Distribution Properties]
For a finite-state Markov chain with transition matrix $P$:
\begin{itemize}
  \item A stationary distribution always exists.
  \item If the chain is irreducible, the stationary distribution is unique.
  \item If the chain is irreducible and aperiodic, the distribution of $X_n$ converges to the stationary distribution as $n\to\infty$ (you can cite this without proof).
\end{itemize}
Irreducible means every state can reach every other in finitely many steps. Aperiodicity rules out deterministic cycles such as a two-state alternation.
\end{remark}

\section{L15}

\begin{definition}
  [Irreducibility and Aperiodicity]
A Markov chain with transition matrix $P$ is \emph{irreducible} if for every pair of states $x_i,x_j\in S$, there exists $k\ge 0$ such that $(P^k)(i,j) > 0$. The chain is \emph{aperiodic} if for every state $x_i$, the greatest common divisor of $\{k\ge 1 : (P^k)(i,i) > 0\}$ is 1.
\end{definition}

\begin{definition}
  [Period]
The \emph{period} of a state $x_i$ is the greatest common divisor of $\{k\ge 1 : (P^k)(i,i) > 0\}$. If the period is greater than 1, the chain exhibits deterministic cycles and does not converge to a stationary distribution.
\end{definition}

\begin{theorem}
  [Markov Chain Convergence]
If a Markov chain with transition matrix $P$ is irreducible and aperiodic, then there exists a unique stationary distribution $\pi$, and for any initial distribution $\mu$,
\[
\lim_{n\to\infty} \mu P^n = \pi.
\]
Furthermore, the convergence is geometric in total variation distance, i.e.\ there exist constants $C>0$ and $\rho\in(0,1)$ such that for all $n$,
\[
\|\mu P^n - \pi\|_{\mathrm{TV}} \le C \rho^n.
\]
Where $\|\cdot\|_{\mathrm{TV}}$ denotes the total variation distance between probability distributions.
\[
\|p-q\|_{\mathrm{TV}} = \frac{1}{2} \sum_{x\in S} |p(x) - q(x)|
\]
\end{theorem}

\section{L16}

\begin{definition}
  [Reversible Markov Chain]
A Markov chain with transition matrix $P$ is \emph{reversible} with respect to a distribution $\pi$ if $\pi$ satisfies the detailed balance equations
\[
\pi(i)P(i,j) = \pi(j)P(j,i)\quad \text{for all } i,j.
\]
\end{definition}

\begin{definition}
  [Spectral Decomposition of Reversible Markov Chains]
If a Markov chain is reversible with respect to $\pi$, then $P$ is self-adjoint on the Hilbert space $L^2(\pi)$, and admits an orthonormal basis of eigenvectors $\{f_k\}$ with corresponding real eigenvalues $\{\lambda_k\}$. The largest eigenvalue is $\lambda_1 = 1$ with eigenvector $f_1 \equiv 1$, and all other eigenvalues satisfy $|\lambda_k| < 1$.
\end{definition}







\end{document}